% ============================================================
% CV - KOSSI SAMUEL GABIAM
% Version correspondant au PDF original
% Compilation : pdfLaTeX
% ============================================================

\documentclass[10pt,a4paper]{article}

% ============================================================
% PACKAGES
% ============================================================

\usepackage[T1]{fontenc}
\usepackage[utf8]{inputenc}
\usepackage[french]{babel}
\usepackage{mathptmx}

\usepackage[
    top=0.9cm,
    bottom=0.9cm,
    left=1.0cm,
    right=1.0cm
]{geometry}

\usepackage{hyperref}
\usepackage{enumitem}
\usepackage{titlesec}
\usepackage{tabularx}
\usepackage{array}
\usepackage{xcolor}
\usepackage{fancyhdr}
\usepackage{needspace}

% ============================================================
% COULEURS
% ============================================================

\definecolor{sectioncolor}{HTML}{1E40AF}
\definecolor{secondary}{HTML}{475569}
\definecolor{linkcolor}{HTML}{1E40AF}

% ============================================================
% HYPERLIENS
% ============================================================

\hypersetup{
    colorlinks=true,
    urlcolor=linkcolor,
    linkcolor=linkcolor,
    pdfborder={0 0 0}
}

% ============================================================
% PAGE
% ============================================================

\pagestyle{fancy}
\fancyhf{}
\renewcommand{\headrulewidth}{0pt}
\fancyfoot[C]{\small\color{secondary}\thepage}

% ============================================================
% SECTIONS
% ============================================================

\titleformat
    {\section}
    {\normalsize\bfseries\color{sectioncolor}}
    {}
    {0pt}
    {}
    [\vspace{-4pt}\color{sectioncolor}\rule{\linewidth}{0.8pt}]

\titlespacing*{\section}
    {0pt}
    {7pt}
    {4pt}

% ============================================================
% LISTES
% ============================================================

\setlist[itemize]{
    leftmargin=1.25em,
    itemsep=1.5pt,
    topsep=1pt,
    parsep=0pt,
    partopsep=0pt
}

% ============================================================
% MACRO POUR LES ENTRÉES
% ============================================================

\newcommand{\cventry}[4]{%
    \needspace{4\baselineskip}%
    \noindent
    \begin{tabularx}{\linewidth}{
        @{}
        X
        >{\raggedleft\arraybackslash}p{3.2cm}
        @{}
    }
        \textbf{\large #1}
        &
        {\small\color{secondary}\textit{#2}}
    \end{tabularx}
    \par
    \noindent\textbf{\textit{#3}}
    \par
    \noindent{\small #4}
    \par
    \vspace{2.5pt}
}

% ============================================================
% DOCUMENT
% ============================================================

\begin{document}

% ============================================================
% EN-TÊTE
% ============================================================

\begin{center}

{\LARGE\bfseries\scshape KOSSI SAMUEL GABIAM}

\par
\vspace{3pt}

{\normalsize\bfseries
Systèmes, Réseaux \& Cloud
\textbullet{}
Développeur Full-Stack}

\par
\vspace{5pt}

{\small
\href{mailto:gabiam.k.samuel@gmail.com}
{gabiam.k.samuel@gmail.com}
\quad
\textbullet
\quad
+228 90 84 77 68
\quad
\textbullet
\quad
Lomé, Togo
}

\par
\vspace{2pt}

{\small
\href{https://www.linkedin.com/in/samuel-g-107663165/}
{linkedin.com/in/samuel-g-107663165}
\quad
\textbullet
\quad
\href{https://github.com/Shieldish}
{github.com/Shieldish}
\quad
\textbullet
\quad
\href{https://gitlab.com/Shieldish}
{gitlab.com/Shieldish}
}

\par
\vspace{2pt}

{\small
\href{https://portfolio.gabiam-k-samuel.workers.dev/}
{portfolio.gabiam-k-samuel.workers.dev}
}

\end{center}

\vspace{-2pt}

\noindent
{\color{sectioncolor}\rule{\linewidth}{1pt}}

\vspace{2pt}

% ============================================================
% PROFIL PROFESSIONNEL
% ============================================================

\section{Profil Professionnel}

Titulaire d'un \textbf{Master Professionnel en Systèmes, Réseaux \& Cloud Computing}
et d'une \textbf{Licence en Génie Logiciel}, je suis développeur Full-Stack
spécialisé dans la conception et l'intégration de systèmes informatiques,
d'applications web et mobiles et d'infrastructures distribuées.

Mon expérience couvre notamment les architectures distribuées, les API,
les bases de données, la sécurité applicative, l'authentification,
le Cloud, la conteneurisation et les pipelines DevOps. Je possède une
approche orientée résolution de problèmes, fiabilité des systèmes et
amélioration des performances, ainsi qu'une expérience dans le développement
de solutions utilisées en production au Togo.

% ============================================================
% EXPÉRIENCE PROFESSIONNELLE
% ============================================================

\section{Expérience Professionnelle}

\cventry
{Développeur Full-Stack}
{Juin 2025 -- Présent}
{SuiSco Sarl}
{Lomé, Togo}

\begin{itemize}

    \item Conception et développement de solutions web et de services
    numériques destinés à différents usages et canaux, avec une attention
    particulière portée à la fiabilité, la sécurité et la performance.

    \item \textbf{EduBoost :} architecture et développement d'une plateforme
    nationale multi-canaux intégrant Web, SMS, USSD et Mobile Money,
    actuellement en production au Togo.

    \item Conception d'une architecture distribuée composée de
    \textbf{15 nœuds Django d'écriture} et d'un nœud central d'agrégation,
    avec mécanisme de génération de tickets anti-collision basé sur PostgreSQL.

    \item Mise en œuvre de traitements asynchrones et distribués avec
    \textbf{Celery, Redis et PostgreSQL}, incluant synchronisation,
    reprise après incident et traitement des demandes Web.

    \item Intégration de mécanismes de sécurité applicative :
    \textbf{OTP, HMAC-SHA256} pour la vérification des webhooks,
    rate-limiting Redis et authentification centralisée.

    \item Développement et maintenance de plusieurs solutions complémentaires :
    \textbf{DefiQuiz, SuiSco My Account} et intégration du
    \textbf{Single Sign-On Keycloak} basé sur OAuth2/OIDC.

\end{itemize}

% ============================================================
% PROJETS MAJEURS
% ============================================================

\section{Projets Majeurs}

\cventry
{DefiQuiz -- Plateforme de Quiz Éducatif}
{2025 -- Présent}
{SuiSco Sarl / YAS TOGO}
{Plateforme Web}

\begin{itemize}

    \item Développement d'une plateforme de quiz multi-clients
    proposant plus de \textbf{5\,000 questions}.

    \item Conception du frontend et du backend avec
    \textbf{Nuxt.js / Next.js, Hono.js et Deno.js}.

\end{itemize}

\cventry
{SuiSco My Account \& SSO Centralisé}
{2025}
{SuiSco Sarl}
{Applications internes}

\begin{itemize}

    \item Développement d'une application web multi-clients de gestion
    de comptes en \textbf{PHP Laravel, API REST et MySQL}.

    \item Mise en place d'un système de \textbf{Single Sign-On}
    avec Keycloak et OAuth2/OIDC pour centraliser et sécuriser
    l'authentification.

\end{itemize}

\cventry
{Plateforme de Gestion de Stages}
{2024}
{Djagora}
{Projet de Fin d'Études}

\begin{itemize}

    \item Développement d'une application Web et Mobile pour la gestion
    des offres de stage, des candidatures et du suivi des étudiants.

    \item Technologies :
    \textbf{React, Node.js, MySQL et React Native}.

\end{itemize}

\cventry
{Extraction \& Visualisation de Données Sociales}
{2021}
{DK-Soft}
{Projet de Stage}

\begin{itemize}

    \item Développement d'une application d'extraction et de visualisation
    analytique en temps réel de données issues des réseaux sociaux.

    \item Mise en œuvre de la \textbf{Stack ELK} :
    Elasticsearch, Logstash et Kibana, avec Angular.

\end{itemize}

% ============================================================
% FORMATION ACADÉMIQUE
% ============================================================

\section{Formation Académique}

\cventry
{Master Professionnel en Systèmes, Réseaux \& Cloud Computing}
{2022 -- 2024}
{Faculté des Sciences de Sfax}
{Sfax, Tunisie}

\begin{itemize}

    \item Ingénierie des systèmes informatiques, architectures distribuées,
    virtualisation, Cloud Computing, administration Linux et sécurité réseau.

\end{itemize}

\cventry
{Licence Nationale en Science de l'Informatique}
{2019 -- 2022}
{Faculté des Sciences de Sfax}
{Sfax, Tunisie}

\begin{itemize}

    \item Spécialité \textbf{Génie Logiciel \& Systèmes d'Information}.

    \item Algorithmique, bases de données, UML, génie logiciel
    et réseaux informatiques.

\end{itemize}

\cventry
{Baccalauréat Scientifique -- Série D}
{2014 -- 2017}
{Lycée La Grâce}
{Sokodé, Togo}

% ============================================================
% COMPÉTENCES TECHNIQUES
% ============================================================

\section{Compétences Techniques}

\begin{tabularx}{\linewidth}{
    @{}
    >{\bfseries\raggedright\arraybackslash}p{3.0cm}
    @{\hspace{8pt}}
    >{\raggedright\arraybackslash}X
    @{}
}

Développement &
React, Next.js, Vue.js, Nuxt 3, Angular,
TypeScript, Java, Go, Django, DRF, Node.js,
PHP Laravel, Hono.js, Deno.js, Python,
API REST, API GraphQL
\\[3pt]

Bases de données &
PostgreSQL, MySQL, Redis, MongoDB,
Elasticsearch
\\[3pt]

Systèmes \& Cloud &
Linux/Unix, Docker, Docker Compose,
architectures distribuées, virtualisation,
Keycloak, SSO
\\[3pt]

DevOps \& Méthodes &
Git, GitHub, GitLab, CI/CD, Celery,
Gunicorn, Agile / Scrum
\\[3pt]

Sécurité &
OAuth2/OIDC, HMAC-SHA256, OTP,
rate-limiting, authentification et sécurisation
des API
\\[3pt]

Intelligence artificielle &
Outils d'IA générative pour le développement,
la recherche, l'analyse, la documentation et
la résolution de problèmes.
Capacité à concevoir, entraîner et
expérimenter des modèles d'IA.
\\[3pt]

Mobile \& Data &
React Native, Flutter, ELK, Elasticsearch,
Kibana, Loki, Grafana, Prometheus

\end{tabularx}

% ============================================================
% LANGUES
% ============================================================

\vspace{4pt}

\noindent
\textbf{Langues :}
Français (langue maternelle)
\quad
\textbullet
\quad
Anglais (compétence professionnelle)

% ============================================================
% FIN
% ============================================================

\end{document}